\section{TIM-MARS Algorithm}

Algorithm~\ref{alg:tim_mars} summarises the implementation-level update used by TIM-MARS.

\begin{algorithm}[t]
\caption{TIM-MARS selected-target memory update}
\label{alg:tim_mars}
\begin{algorithmic}[1]
\REQUIRE tracker candidates $\mathcal{C}_t$, previous memory $m_{t-1}$
\ENSURE output target $y_t$, state $s_t$
\STATE discard candidates with $q_t^i < \tau_q$
\IF{$\mathcal{C}_t' = \emptyset$}
    \STATE update missing-frame counter
    \STATE $s_t \leftarrow \mathrm{UNCERTAIN}$ or $\mathrm{LOST}$
    \STATE $y_t \leftarrow \emptyset$
    \RETURN
\ENDIF
\FORALL{$c_t^i \in \mathcal{C}_t'$}
    \STATE compute $I_i$, $D_i$, $Z_i$, $Q_i$, and $B_i$
    \STATE compute base score $S_i^0$
\ENDFOR
\STATE compute base ambiguity margin $\Delta_t^0$
\STATE decide whether appearance should be used
\FORALL{$c_t^i \in \mathcal{C}_t'$}
    \IF{appearance is available and geometry allows appearance}
        \STATE compute raw MARS similarity $A_i$
        \IF{appearance is enabled and $A_i \geq \tau_A$}
            \STATE $S_i \leftarrow \mathrm{clip}(S_i^0+w_AA_i,0,1)$
        \ELSE
            \STATE $S_i \leftarrow S_i^0$
        \ENDIF
    \ELSE
        \STATE $S_i \leftarrow S_i^0$
    \ENDIF
\ENDFOR
\STATE $c_t^\star \leftarrow \arg\max_{c_t^i \in \mathcal{C}_t'} S_i$
\STATE choose threshold $\tau_s$ from previous state
\IF{$id(c_t^\star) = id_{t-1}^\star$}
    \STATE relax threshold by $\delta_{id}$
\ENDIF
\IF{$S_t^\star < \tau_s$}
    \STATE $s_t \leftarrow \mathrm{UNCERTAIN}$ or $\mathrm{LOST}$
    \STATE $y_t \leftarrow \emptyset$
    \RETURN
\ENDIF
\IF{conservative appearance filtering is enabled}
    \IF{appearance was not used and appearance is required}
        \STATE $s_t \leftarrow \mathrm{UNCERTAIN}$ or $\mathrm{LOST}$
        \STATE $y_t \leftarrow \emptyset$
        \RETURN
    \ENDIF
    \IF{appearance was used and $(A^\star < \tau_C$ or $\Delta_A < \tau_M)$}
        \STATE $s_t \leftarrow \mathrm{UNCERTAIN}$ or $\mathrm{LOST}$
        \STATE $y_t \leftarrow \emptyset$
        \RETURN
    \ENDIF
\ENDIF
\IF{previous state was $\mathrm{UNCERTAIN}$ or $\mathrm{LOST}$}
    \STATE $s_t \leftarrow \mathrm{REACQUIRED}$
\ELSE
    \STATE $s_t \leftarrow \mathrm{LOCKED}$
\ENDIF
\STATE update selected-target memory
\STATE update appearance memory only in confirmed LOCKED state
\STATE $y_t \leftarrow b_t^\star$
\end{algorithmic}
\end{algorithm}

Figure~\ref{fig:state_machine} shows the state-machine interpretation.

\begin{figure}[t]
\centering
\begin{tikzpicture}[
    node distance=8mm,
    state/.style={draw, rounded corners, align=center, minimum width=20mm, minimum height=7mm, font=\footnotesize},
    arrow/.style={-Latex, thick, font=\scriptsize}
]
\node[state] (no) {NO\\TARGET};
\node[state, right=of no] (locked) {LOCKED};
\node[state, below=of locked] (uncertain) {UNCERTAIN};
\node[state, left=of uncertain] (lost) {LOST};
\node[state, right=of uncertain] (reacq) {REACQUIRED};

\draw[arrow] (no) -- node[above] {select} (locked);
\draw[arrow] (locked) -- node[right] {weak or ambiguous} (uncertain);
\draw[arrow] (uncertain) -- node[above] {timeout} (lost);
\draw[arrow] (lost) -- node[below] {candidate accepted} (reacq);
\draw[arrow] (reacq) -- node[right] {confirmed} (locked);
\draw[arrow] (locked) edge[loop above] node {accepted} (locked);
\draw[arrow] (lost) edge[loop left] node {no safe candidate} (lost);
\end{tikzpicture}
\caption{TIM-MARS state-machine interpretation. Ambiguous candidates are rejected into UNCERTAIN or LOST instead of being published as control-valid targets.}
\label{fig:state_machine}
\end{figure}
